Debate has been recently proposed as a training protocol that is hoped to
produce more aligned models even as model capabilities scale. The idea is
that the human, instead of providing supervision in the form of labels,
judges a debate between two models that compete to prove the correctness of
opposite labels. It can be shown that the set of functions that can be
correctly judged in this way is a lot more complex than the set of functions
that can be directly evaluated, thus suggesting that perhaps debate can be
used to teach very complicated functions to models. However, previous work
has largely ignored the fact that this is a \emph{training} protocol, and
simply analyzes the game being played between the two debaters for a known
judge. We take the first steps towards understanding the learning problem
implicit within the protocol by formulating it as an online learning
problem. We show that once we formulate the learning problem precisely,
there are cases when the best way to learn given access to a judge is
different from running an adversarial debate. This motivates us to define a
cooperative version of the protocol, whose learnability we characterize
using a combinatorial dimension. We also consider the possibility of the
human judge being fallible and occasionally judging the debate incorrectly.
We show that a small number of errors in judgement can lead to a large error
in learning.
